\begin{contribucion}{Rafael López Rodríguez}

\begin{itemize}

	\item Investigué sobre las motivaciones sociales y técnicas detrás de la creación de deepfakes, así como los problemas críticos que presenta la detección de audio clonado en la actualidad.
	
	\item Investigué sobre el estado del arte de los modelos de clasificación Regresión Logística y K-Nearest Neighbors (k-NN), profundizando en sus fundamentos matemáticos, sus hiperparámetros y su idoneidad para el manejo de características acústicas.
	
	\item Investigué exhaustiva sobre las técnicas de IA Explicable (XAI), analizando cuáles de ellas son aplicables específicamente al dominio del audio y la detección de fraude. Esto incluyó tanto técnicas post-hoc (que explican modelos ya entrenados) como técnicas ante-hoc (modelos interpretables por diseño).
	
	\item Estudié la viabilidad de aplicar arquitecturas intrínsecamente explicables como ProtoPNet y SincNet para la detección de audio, realizando pruebas para determinar si estas técnicas de explicabilidad ante-hoc aportarían valor frente a los métodos tradicionales.
	
	\item Desarrollé el código necesario para la implementación de los algoritmos de Regresión Logística y k-NN. Diseñé un flujo de trabajo que incluyó el ajuste fino de hiperparámetros y un análisis detallado de las métricas de rendimiento para optimizar la capacidad de detección.
	
	\item Apliqué técnicas de explicabilidad post-hoc a los modelos desarrollados, implementando SHAP (global y local), LIME (local) y DiCE (generación de explicaciones contrafactuales) para interpretar las decisiones del sistema.
	
	\item Redacté los apartados de Objetivos y Planificación dentro de la Introducción del proyecto.
	
	\item Redacté, en el bloque de Estado del Arte, las secciones correspondientes a los modelos de clasificación (Regresión Logística y k-NN) y a la totalidad de las técnicas de explicabilidad analizadas (SHAP, LIME, DiCE, Integrated Gradients, Grad-CAM, SincNet y ProtoPNet).
	
	\item Elaboré los apartados técnicos de la Metodología relativos a la implementación y el análisis de resultados de los modelos de Regresión Logística y k-NN.

\end{itemize}
\end{contribucion}
